\documentclass[11pt,a4paper]{article}
\usepackage[T1]{fontenc}
\usepackage{lmodern}
\usepackage[margin=26mm]{geometry}
\usepackage{amsmath,amssymb,amsthm,booktabs,array,longtable}
\usepackage{microtype,multicol}
\usepackage[hidelinks]{hyperref}
\setlength{\parskip}{3pt}
\setlength{\emergencystretch}{2em}
\raggedbottom
\allowdisplaybreaks[1]
\clubpenalty=10000
\widowpenalty=10000
\newtheorem{theorem}{Theorem}[section]
\newtheorem{lemma}[theorem]{Lemma}
\newtheorem{corollary}[theorem]{Corollary}
\theoremstyle{remark}
\newtheorem{remark}[theorem]{Remark}
\newcommand{\Z}{\mathbb Z}
\newcommand{\Q}{\mathbb Q}
\newcommand{\N}{\mathbb N}
\newcommand{\F}{\mathbb F}
\newcommand{\ps}[1]{#1[[x]]}
\newcommand{\Rev}{\operatorname{Rev}}
\newcommand{\rev}{^{\langle-1\rangle}}
\newcommand{\Plus}{\operatorname{Plus}}
\newcommand{\Minus}{\operatorname{Minus}}
\newcommand{\coeff}[2]{[x^{#1}]#2}
\newcommand{\E}{\mathcal E}
\Urlmuskip=0mu plus 1mu
\hypersetup{pdftitle={A380680: A false square-index conjecture and its twice-square correction},pdfauthor={Dev Goyal},pdfsubject={Separate manuscript with AI disclosure and Lean proof map}}
\title{A380680\\[4pt]\large A false square-index conjecture and its twice-square correction}
\author{Dev Goyal}
\date{24 October 2026}
\begin{document}
\maketitle
\begin{center}
\fbox{\begin{minipage}{0.9\linewidth}\small
\textbf{AI disclosure.} ChatGPT (OpenAI) generated the mathematical arguments,
Lean proofs, and this manuscript in response to Dev Goyal's requests.
Dev Goyal is the sole human author. The source sequence and conjecture or
observation are credited to Paul D. Hanna.
\textbf{Independent human review is pending.}
\end{minipage}}
\end{center}
\begin{abstract}
The retrieved square-index conjecture fails already at $a_1=4$ and $a_2=6$. For the explicit positive-index rational theta sum, we prove the corrected identity $B\equiv1+2\sum_{j\ge1}x^{2j^2}\pmod4$. The full exceptional support is therefore twice a square, and both clauses of the false statement are explicitly refuted.
\end{abstract}

\section{Source, notation, and scope}

The source is Paul D. Hanna's OEIS entry A380680 \cite{oeis680},
retrieved on 24 October 2026. Its initial sequence index is 0.
The constant coefficient is one; the residue classification below
is restricted to positive indices.

All series are formal. The set $\N$ includes zero, and $[x^n]U$ denotes a
coefficient. A normalized series is $x+O(x^2)$; $O(x^r)$ means divisibility
by $x^r$, over any coefficient ring used below. We write $U\circ V=U(V(x))$
only when $V(0)=0$. Compositional inversion, denoted $\Rev(U)$ or $U\rev$,
is distinct from multiplicative inversion of a unit. Congruences are
coefficientwise; ordinary products and compositional iterates are distinct.

\section{Units, coefficient maps, and doubling}

A constant-one series $U=1+\sum_{n>0}u_nx^n$ has a unique multiplicative
inverse $V=1+\sum_{n>0}v_nx^n$: the coefficient equation for $UV=1$ is
\[
 v_n=-\sum_{k=1}^n u_kv_{n-k}\qquad(n>0).
\]
Thus inversion uses only ring operations, is defined integrally, and commutes
with coefficient homomorphisms. If two constant-one series agree modulo
$x^r$, the recurrence shows that their inverses agree modulo $x^r$ as well.
Locally finite coefficient sums, products, and integral powers also commute
with coefficient maps.

\begin{lemma}[Doubling]\label{lem:binary}
For formal series $U,V$ over $\Z/4\Z$, $2U=2V$ if and only if their
reductions modulo two agree.
\end{lemma}
\begin{proof}
For integers $u,v$, $4\mid2(u-v)$ is equivalent to $2\mid u-v$. Apply
this equivalence to every coefficient. No cancellation by two inside
$\Z/4\Z$ is used.
\end{proof}

\section{The literal positive-index source}

For a formal series $U$ over a commutative ring, put
\[
 Q_j(U)=\frac{U-x^j}{1-x^jU}\quad(j\ge1),\qquad
 \Sigma(U)=\sum_{j\ge1}x^{j^2}Q_j(U)^j. \tag{1}
\]
The denominator is a constant-one unit, so the quotient denotes a genuine
multiplicative inverse. The $j$th summand has order at least $j^2$.
Each coefficient is therefore a finite sum, and no analytic convergence
assumption is required. This is the positive-index formula (2) given by OEIS.
A separate formal Laurent-series equivalence with the bilateral title is
outside the claimed Lean coverage.

The selected explicit series is
\[
 B=1+2\Sigma(2). \tag{2}
\]
Here $2$ denotes the constant series. This distinction from inserting
$B$ in the argument of $\Sigma$ determines the different residue support.
\begin{theorem}[The corrected A380680 statement]\label{thm:theta680}
The explicit integral series (2) satisfies
\[
 B\equiv1+2\sum_{j\ge1}x^{2j^2}\pmod4.
\]
Therefore, for every $n>0$, its coefficient is two modulo four if and
only if $n=2j^2$ for some positive integer $j$, and is divisible by four
at all other positive indices.
\end{theorem}

\begin{proof}
For the constant series two, reduction modulo two gives
\[
 Q_j(2)=\frac{2-x^j}{1-2x^j}\equiv x^j\pmod2.
\]
Thus the $j$th summand of $\Sigma(2)$ reduces to
$x^{j^2}(x^j)^j=x^{2j^2}$. Apply the doubling statement in
Lemma~\ref{lem:binary} and then (2). Distinct positive $j$ yield distinct
exponents $2j^2$.
\end{proof}

\begin{corollary}[The retrieved square-index conjecture is false]
For A380680, $[x]B=4$ and $[x^2]B=6$. The retrieved claim that the
residue-two support consists of squares is false, as is its separate
claim of divisibility by four at every positive nonsquare index.
\end{corollary}

\begin{proof}
In degrees one and two, only $j=1$ can contribute to (2). Since
\[
 \frac{2-x}{1-2x}=2+3x+O(x^2),
\]
we obtain $B=1+4x+6x^2+O(x^3)$. The square index one has residue zero,
while the nonsquare index two has residue two. Both failures are proved
as explicit negations of the corresponding universal statements in Lean.
\end{proof}


The displayed source coefficients already expose both failures. The
newly presented theorem supplies the complete corrected twice-square pattern;
it is a correction, not a proof of the false square-index statement.

\section{Separate Lean proof and checked scope}

The entry file is \nolinkurl{lean/Entries/A380680.lean}. Its explicit
theorem statements and proof terms use the namespace
\nolinkurl{FinalCollection.A380680}. The complete imported implementations
and their shared mathematical dependencies are included in the collection.
The following local theorem names give the proof map:

\begin{multicols}{2}\raggedright\begin{itemize}\setlength{\itemsep}{2pt}\setlength{\parsep}{0pt}
\item \nolinkurl{explicit_actual_source}.
\item \nolinkurl{a380680_corrected_statement}.
\item \nolinkurl{a380680_counterexamples}.
\item \nolinkurl{a380680_literal_statement_is_false}.
\item \nolinkurl{a380680_complement_is_false}.
\item \nolinkurl{complete_series_mod_four}.
\end{itemize}\end{multicols}

The main implementation module is \nolinkurl{Round4.ThetaRational}.
The catalogue records the full original endpoint names and the new entry
theorems. The project pins Lean 4.19.0 and Mathlib revision
\nolinkurl{c44e0c8ee63ca166450922a373c7409c5d26b00b}. The collection
report records source recompilation and transitive axiom checks. Accepted
axioms are propositional extensionality, classical choice, and quotient
soundness. All index conclusions above are unbounded; exact finite
coefficient checks in the package provide separate corroboration.

\paragraph{Limits and status.}
The formal source is the explicit positive-index formula (2) in OEIS. A separate formal Laurent-series equivalence to the bilateral title is not claimed. No convergence or real special-value assertion is proved.
The manuscript consolidates completed results; it does not assert global
novelty, exclusive priority, peer review, OEIS acceptance, or a new submission.
Standard algebraic ingredients are not claimed as new discoveries.

\begin{thebibliography}{9}
\raggedright
\bibitem{oeis680}
Paul D. Hanna, \emph{OEIS A380680}, 6 February 2025.
\url{https://oeis.org/A380680}. Retrieved 24 October 2026.
\end{thebibliography}
\end{document}
